\scanpage{13}
\mspage{12.}

d'ordre sur $\Phi$ est discrète, i.e. $\Phi\neq\varnothing$ et formé d'éléments minimaux de $\mathcal{F}$.

Ce sont donc les \deleted{ens. configurations} parties de $\mathcal{F}_{\unsure{e\bar l|\overline{el}|\text{él}}\,\mathrm{min}}$ qui ont les propriétés suivantes

a) c'est une configuration, i.e. \inserted{formé de fig. él.} \unsure{mutuellement|mutuellemt} \deleted{dis} compatibles, et elle est non vide

b) Elle est loc. finie.

Dans les exemples \unsure{types|type}, ceci correspond aux parties \inserted{fermées} $Y$ de $X$ qui sont des sous-variétés \inserted{non vides}, découpées en strates par les composantes connexes.

Aussi, \unsure{\deleted{les}|les} figures de dim $\leq 0$ s'appeleront les \uline{sous}-\uline{variétés} (au sens strict, i.e. avec stratification triviale).

La sous-variété vide correspond : la configuration vide.
